\section{Rank spreading and the surviving order}
\label{sec:ranks}

The distribution of rank embeddings used in the construction makes comparisons
at different triples of positions nearly indistinguishable.  We first prove
this statement, and then show how it turns a surviving topological order into
compatible monotone Boolean functions.

\subsection{A distribution on increasing embeddings}

For a finite set $S$, write $\binom{S}{3}$ for its three-element subsets,
each represented in increasing order when $S$ is ordered.  For probability
measures $\mu,\mu'$ on a common finite space $\Omega$, we use the convention
\[
  \TV(\mu,\mu')=\max_{\mathcal A\subseteq\Omega}
     \abs{\mu(\mathcal A)-\mu'(\mathcal A)}.
\]

\begin{lemma}[Rank spreading]
\label{lem:rank-spreading}
For every integer $L\ge3$ and every $\eta>0$, there are an integer $N\ge L$
and a rational probability distribution on increasing maps
$\psi\colon[L]\hookrightarrow[N]$ such that, for all
$I,I'\in\binom{[L]}{3}$,
\begin{equation}
\label{eq:rank-spreading-tv}
  \TV\bigl(\mathcal L(\psi(I)),\mathcal L(\psi(I'))\bigr)\le\eta.
\end{equation}
Here $\mathcal L$ denotes the law of a random variable, and all the laws in
Equation~\eqref{eq:rank-spreading-tv} use the same random map $\psi$.
\end{lemma}

\begin{proof}
Choose a positive number $\delta<\eta/2$.  Let
$\mathcal I=\binom{[L]}{3}$.  Partition $[0,1]$ into a finite number $b$ of
intervals, each of diameter at most $\delta$, assigning boundary points to one
of their adjacent intervals.  The finite Ramsey theorem for triples
\cite[Theorem~B]{Ramsey1930} gives an $N$ such that every coloring of
$\binom{[N]}{3}$ with at most $b^{\abs{\mathcal I}^2}$ colors has a
monochromatic set of size $L$.

Consider the following finite zero-sum game.  The minimizing player chooses
an increasing map $\psi\colon[L]\hookrightarrow[N]$.  The maximizing player
chooses a triple $(I,I',\mathcal A)$, where $I,I'\in\mathcal I$ and
$\mathcal A\subseteq\binom{[N]}{3}$.  The payoff to the maximizing player is
\[
  \ind{\psi(I)\in\mathcal A}-\ind{\psi(I')\in\mathcal A}.
\]
Fix an arbitrary mixed strategy $\tau$ of the maximizing player.  Let
$w_{I,I'}$ be its marginal probability of the ordered pair $(I,I')$.  When
$w_{I,I'}>0$, define, for $Z\in\binom{[N]}{3}$,
\[
  q_{I,I'}(Z)=\Pr_{\tau}[Z\in\mathcal A\mid I,I'];
\]
put $q_{I,I'}(Z)=0$ when $w_{I,I'}=0$.
Color $Z$ by the vector of intervals containing $q_{I,I'}(Z)$, one interval
for every ordered pair $(I,I')\in\mathcal I^2$.  The number of possible
colors is at most $b^{\abs{\mathcal I}^2}$, so there is a monochromatic
$L$-element set.  Let $\psi$ be its increasing enumeration.  For every
ordered pair $(I,I')$, both $\psi(I)$ and $\psi(I')$ are triples in this
set, and hence
\[
  \abs{q_{I,I'}(\psi(I))-q_{I,I'}(\psi(I'))}\le\delta.
\]
The expected payoff against $\tau$ is consequently at most
\[
  \sum_{I,I'\in\mathcal I}w_{I,I'}\,
   \bigl(q_{I,I'}(\psi(I))-q_{I,I'}(\psi(I'))\bigr)
  \le\delta.
\]
The number $N$ was chosen from a bound on the number of colors that does
not depend on $\tau$.  Thus every mixed strategy of the maximizing player
has a pure response with payoff at most $\delta$.

By the finite minimax theorem \cite[Section~II.3]{vonNeumann1928}, there is a probability
distribution $\mu$ on increasing maps for which the expected payoff against
every pure test is at most $\delta$.  Choose a rational probability
distribution $\widehat\mu$ on the same finite space with
\[
  \sum_{\psi}\abs{\widehat\mu(\psi)-\mu(\psi)}<\eta/2.
\]
Such distributions are dense in the finite probability simplex.  Since
each payoff lies in $[-1,1]$, this replacement changes any expected payoff
by less than $\eta/2$.  Thus every pure test has expected payoff less than
$\delta+\eta/2<\eta$ under $\widehat\mu$.  The ordered pairs include
both $(I,I')$ and $(I',I)$, and $\mathcal A$ ranges over all subsets of
$\binom{[N]}{3}$.  The resulting inequalities give
Equation~\eqref{eq:rank-spreading-tv}.
\end{proof}

\subsection{Boolean functions from rank comparisons}

For the remainder of this section and
Sections~\ref{sec:prefix}--\ref{sec:soundness}, suppose that the constructed
weighted digraph has a feedback vertex set $F$ of cost at most $D$.
Fix a total topological order $\prec$ of the surviving vertices.  Every
rank vertex survives, since it has weight $D+1$.  In particular, the
constant vertices satisfy
\[
  o_1\prec o_2\prec\cdots\prec o_N.
\]
All subsequent random data are sampled after $F$ and this order have been
fixed.  In a test family of total mass $W$, the deletion cost paid by $F$
in that family is exactly $W$ times the probability that its sampled test
vertex belongs to $F$.

For an allowed tuple $\mathbf d$, a subset $E\subseteq\Lambda_{\mathbf d}$,
and indices $1\le a<c<h\le L$, define
\begin{equation}
\label{eq:rank-comparison-bits}
  G_{\mathbf d}(E;a,c,h)
   =\ind{b_{\mathbf d}^{\psi}(E;a,h)\prec o_{\psi(c)}},
  \qquad
  g_{\mathbf d}(E)=G_{\mathbf d}(E;2,15,L-1).
\end{equation}
The dependence of these functions on the sampled embedding $\psi$ is
suppressed in the notation.

\begin{definition}[Good tuples]
\label{def:good-tuple}
A triple $(a,c,h)$ is \emph{usable} if
\begin{equation}
\label{eq:usable-triples}
  5\le c\le L-5,\qquad 2\le a<c,\qquad c+2\le h\le L.
\end{equation}
The tuple $\mathbf d$ is \emph{good at $\psi$} if
\[
  G_{\mathbf d}(E;a,c,h)=g_{\mathbf d}(E)
\]
for every $E\subseteq\Lambda_{\mathbf d}$ and every usable triple.
Otherwise it is \emph{bad at $\psi$}.
\end{definition}

\begin{lemma}[Probability of goodness]
\label{lem:goodness}
For each fixed allowed tuple $\mathbf d$,
\begin{equation}
\label{eq:goodness-probability}
  \Pr_{\psi}[\mathbf d\text{ is bad at }\psi]
  \le 3\eta L^3 2^{\abs{\Lambda_{\mathbf d}}}
  \le\rho.
\end{equation}
The same upper bound $\rho$ holds when $\mathbf d$ is sampled independently
of $\psi$.
\end{lemma}

\begin{proof}
Fix $\mathbf d$ and $E\subseteq\Lambda_{\mathbf d}$.  For integers
$r_1<r_2<r_3$ in $[N]$, let $x_{r_1,r_3}$ be the profile taking value
$r_1$ on $E$ and $r_3$ on its complement.  Once the surviving order is
fixed, the function
\[
  (r_1,r_2,r_3)\longmapsto
    \ind{b_{\mathbf d}(x_{r_1,r_3})\prec o_{r_2}}
\]
is a deterministic Boolean function of this triple alone.  Applying
Lemma~\ref{lem:rank-spreading} to it shows that
\begin{equation}
\label{eq:comparison-expectations}
  \abs{\E_{\psi}G_{\mathbf d}(E;a,c,h)
       -\E_{\psi}G_{\mathbf d}(E;a',c',h')}\le\eta
\end{equation}
for any two increasing triples of positions, whether usable or not.

Equation~\eqref{eq:comparison-expectations} compares expectations; it does
not by itself control disagreement under the same embedding.  We obtain
that control through pointwise ordered pairs of bits, for which the
probability of disagreement equals the absolute difference of expectations.
Fix a usable triple $(a,c,h)$ and abbreviate
\[
\begin{split}
  A_0&=G_{\mathbf d}(E;2,15,L-1),\\
  A_1&=G_{\mathbf d}(E;2,c,L-1),\\
  A_2&=G_{\mathbf d}(E;1,c,c+1),\\
  A_3&=G_{\mathbf d}(E;a,c,h).
\end{split}
\]
The bits $A_0$ and $A_1$ compare the same rank vertex to two constant
vertices.  They are pointwise ordered: $A_0\le A_1$ if $c\ge15$, and
$A_1\le A_0$ if $c\le15$.  Equation~\eqref{eq:comparison-expectations}
therefore gives $\Pr[A_0\ne A_1]\le\eta$.

The profile of $b_{\mathbf d}^{\psi}(E;1,c+1)$ is strictly smaller at
every coordinate than the profile of
$b_{\mathbf d}^{\psi}(E;2,L-1)$: on $E$ use $1<2$, and on its
complement use $c+1<L-1$.  The strict rank arc from the former vertex to
the latter forces $A_2\ge A_1$.  Similarly, $1<a$ and $c+1<h$ imply
$A_2\ge A_3$.  Applying Equation~\eqref{eq:comparison-expectations}
to these ordered pairs gives
\[
  \Pr[A_1\ne A_2]\le\eta,
  \qquad \Pr[A_2\ne A_3]\le\eta.
\]
A union bound yields
\[
  \Pr\bigl[G_{\mathbf d}(E;a,c,h)\ne g_{\mathbf d}(E)\bigr]
  \le3\eta.
\]
There are at most $L^3$ usable triples and
$2^{\abs{\Lambda_{\mathbf d}}}$ subsets $E$.  A second union bound proves
the first inequality of Equation~\eqref{eq:goodness-probability}.
Every allowed tuple has length at most $T$ and each slot alphabet has
size at most $m$, so $\abs{\Lambda_{\mathbf d}}\le m^T$.  The choice
$3\eta L^3 2^{m^T}\le\rho$ proves the second inequality.  Finally,
averaging the pointwise bound over any tuple distribution independent of
$\psi$ gives the last assertion.
\end{proof}

The order in Lemma~\ref{lem:goodness} may depend on the entire weighted
construction, including its rank distribution.  The argument only uses that
the order is fixed before the fresh draw of $\psi$.  Nor does
Equation~\eqref{eq:goodness-probability} assert simultaneous goodness of
all tuples; the subsequent estimates use its stated bound for sampled
tuples.

\begin{lemma}[Monotonicity]
\label{lem:monotonicity}
If $\mathbf d$ is good at $\psi$, then
$g_{\mathbf d}\colon2^{\Lambda_{\mathbf d}}\to\{0,1\}$ is monotone:
\[
  E\subseteq E'\quad\Longrightarrow\quad
  g_{\mathbf d}(E)\le g_{\mathbf d}(E').
\]
Moreover,
\[
  g_{\mathbf d}(\varnothing)=0,
  \qquad g_{\mathbf d}(\Lambda_{\mathbf d})=1.
\]
\end{lemma}

\begin{proof}
For $E\subseteq E'$, compare
\[
  b_{\mathbf d}^{\psi}(E';2,17)
  \quad\text{and}\quad
  b_{\mathbf d}^{\psi}(E;3,18).
\]
The first profile is strictly smaller everywhere: the index pairs in
$E$, $E'\setminus E$, and $\Lambda_{\mathbf d}\setminus E'$ are,
respectively, $(2,3)$, $(2,18)$, and $(17,18)$.  The rank arc consequently
implies
\[
  G_{\mathbf d}(E';2,15,17)
  \ge G_{\mathbf d}(E;3,15,18).
\]
Both triples are usable, so goodness identifies these two bits with
$g_{\mathbf d}(E')$ and $g_{\mathbf d}(E)$.  For the extreme sets, the
reference profile is $o_{\psi(L-1)}$ when $E=\varnothing$ and
$o_{\psi(2)}$ when $E=\Lambda_{\mathbf d}$.  Their positions relative to
$o_{\psi(15)}$ give the claimed values.
\end{proof}

\begin{lemma}[Prefix compatibility]
\label{lem:compatibility}
Let $\mathbf s w$ be an allowed extension of $\mathbf s$.  For every
$E\subseteq\Lambda_{\mathbf s}$ and every increasing triple $(a,c,h)$,
\begin{equation}
\label{eq:prefix-lift-bits}
  G_{\mathbf s w}(E\times\Lambda_w;a,c,h)
  =G_{\mathbf s}(E;a,c,h).
\end{equation}
In particular,
$g_{\mathbf s w}(E\times\Lambda_w)=g_{\mathbf s}(E)$.
If $\mathbf s w$ is good at $\psi$, then so is $\mathbf s$.
The same conclusions hold for a prefix obtained by dropping several final
slots.
\end{lemma}

\begin{proof}
The profile on $\Lambda_{\mathbf s w}$ that takes value $\psi(a)$ on
$E\times\Lambda_w$ and $\psi(h)$ elsewhere is the pullback of the
corresponding profile on $\Lambda_{\mathbf s}$ under the projection
dropping the last slot.  Definition~\ref{def:rank-graph} identifies the
two rank vertices.  Their comparisons to the same constant marker are
identical, proving Equation~\eqref{eq:prefix-lift-bits}.  Apply this
identity both to the reference triple and to every usable triple: goodness
of the extension gives goodness of the shorter tuple for each subset $E$.
Iteration proves the final assertion.  The projection identity holds on
the formal product alphabets even when a game vertex occurs in more than
one slot.
\end{proof}

\subsection{When a wildcard test must be deleted}

\begin{lemma}[Wildcard forcing]
\label{lem:wildcard-forcing}
Consider a wildcard test vertex $v_{\mathrm{test}}$ from
Definition~\ref{def:wildcard}, with tuple $\mathbf d$, embedding $\psi$,
ordered cells $C_1,\ldots,C_k$, and star set $Z$.  Suppose that
$\mathbf d$ is good at $\psi$.  For an integer $d$ with $0\le d\le k$,
put
\[
  E=\Bigl(\bigcup_{\ell=1}^{d}C_\ell\Bigr)\setminus Z.
\]
If
\begin{equation}
\label{eq:wildcard-forcing-criterion}
  g_{\mathbf d}(E)=0,
  \qquad g_{\mathbf d}(E\cup Z)=1,
\end{equation}
then $v_{\mathrm{test}}\in F$.
\end{lemma}

\begin{proof}
Set $c=10d+15$ and define the rank vertices
\[
  v^-=b_{\mathbf d}^{\psi}(E;2,c+2),
  \qquad
  v^+=b_{\mathbf d}^{\psi}(E\cup Z;c-2,L-1).
\]
We first check the wildcard's incident arcs.  If
$\lambda\in C_\ell\setminus Z$ with $\ell\le d$, then
$\lambda\in E$, so the two profiles take values $\psi(2)$ and
$\psi(c-2)$, respectively.  The inequalities
\[
  2<10\ell+10\le10d+10<c-2
\]
place these values strictly below and strictly above
$z_\ell=\psi(10\ell+10)$.  If instead $\ell>d$, then
$\lambda\notin E\cup Z$, so the profiles take values $\psi(c+2)$ and
$\psi(L-1)$.  Here
\[
  c+2=10d+17<10\ell+10
  \le10k_{\max}+10=L-20<L-1,
\]
which gives the required strict inequalities again.  Thus the constructed
digraph contains the arcs
\[
  v^-\longrightarrow v_{\mathrm{test}}\longrightarrow v^+.
\]
Empty cells impose no conditions in these checks.  The first case is
absent when $d=0$, and the second is absent when $d=k$.

Since $0\le d\le k\le k_{\max}$ and $L=10k_{\max}+30$,
\[
  15\le c\le L-15.
\]
Both $(2,c,c+2)$ and $(c-2,c,L-1)$ are therefore usable triples.
Goodness and Equation~\eqref{eq:wildcard-forcing-criterion} show that
$v^-$ does not precede $o_{\psi(c)}$, whereas $v^+$ does precede it.
If $v_{\mathrm{test}}$ survived, the displayed path would give
$v^-\prec v_{\mathrm{test}}\prec v^+\prec o_{\psi(c)}$, a contradiction.
Hence $v_{\mathrm{test}}\in F$.
\end{proof}
